% Pursuing Stacks (cotes 134-2 to 134-8): archive cote 134-4, site page 79 (PDF page 80), typescript page 167.
% typescript. Draft compared: ps3.tex:4043-4111.
% AI-transcribed by Claude Opus 5.5 (claude-opus-5-5, adaptive thinking, effort high), Message Batches API, batch ps2 (msgbatch_01NVwr1Yuo8qJf5rykeGBXPK), 2026-11-12; not yet checked by eye.
\scanpage{134-4/79}
\mspage{167}

\hand{L} b) The functor (2) commutes with finite sums,

possibly even with infinite ones, which should be closely related to
property a), and

\hand{L} c) The functor (2) commutes with finite products,

\hand{(}maybe even with infinite ones, under suitable assumptions\hand{)}.

The following property, due to Quillen for weak equivalences, is
used in order to prove for these (and the cohomological analogs) the
fibration axiom L 5 (what we get directly is the case of cofibrations,
as a matter of fact\xout{\illegible{1}}\hand{ - cf prop. page 97):}

\hand{L} d) If $f$: $C\to C'$ is in $\underline{W}$, so si $f^{o}$: $C^{o}\to C'^{o}$ for the dual categories

\margin{c'est vrai -- pour tout \unsure{localisateur} \unsure{fonda-mental} \unsure{(i.e.} \unsure{un} \underline{W} \unsure{qui satisf.} \unsure{L1, L3, L4)} \unsure{d'après} \unsure{qui vérifie} \unsure{L5, L4.} \illegible{1}}
This is about all I have in mind at present, as far as further
properties of a $\underline{W}$ are concerned. The property c) however brings to \xout{min}
mind the natural \inserted{(weaker)} condition that the cartesian product of two maps in $\underline{W}$
should equally be in $\underline{W}$ . The argument in the beginning of par. 40
(p.69) carries over here, and we get :

\uline{Proposition} Let $\underline{W}$ be a basic modelizer, and $C$ in (Cat) such that \inserted[typed]{\hand{(} $C$ is $\underline{W}$-aspheric i.e.\hand{)}}
$C\to e$ is in $\underline{W}$ , than $\underline{C}$ is even $\underline{W}$-aspheric over $e$, i.e. for any $A$ in
$C$, $C\mathrm{x}A\to A$ is in $\underline{W}$ .

However, the proof given for the more general statement we have in
mind (cf proposition on page 69) does not carry over using only the loca-%
lization axiom in the form L 4) it was stated above, as far as I can
see. This suggests a \xout{more general} stronger version L 4') of L 4) which
we may have to use eventually, relative to a \xout{situation} commutative
triangle in (Cat) as on page 70
\diagram{
\begin{tikzcd}[column sep=small]
P' \ar[rr,"F"] \ar[dr] & & P \ar[dl] \\
& C &
\end{tikzcd}
}
,

when assuming that the induced maps (for arbitrary $c$ in $C$)
\[ F_{/c} : P'_{/c}\to P_{/c} \]
are in $\underline{W}$ , to deduce that $F$ itself is in $\underline{W}$ . However this ``strong
localization axiom\deleted{e}'' is a consequence (as is the weaker one L 4)) of
\margin{$\times$} the fibration axiom \inserted{L 5)}, which implies also directly the property we have
in mind , namely \deleted{\xout{f,g $\underline{W}$ implies}}
\[ f,g \hand{\in} \underline{W} \text{ implies } f\mathrm{x}g \hand{\in} \underline{W} \quad . \]

\margin{79}
